\documentclass[11pt]{article}
\usepackage{amsmath,amssymb}
\usepackage{lmodern}
\usepackage[margin=1in]{geometry}
\usepackage{microtype}
\usepackage[hidelinks]{hyperref}
\linespread{1.05}
\setlength{\parindent}{0pt}
\setlength{\parskip}{4pt plus 1pt minus 1pt}
\postdisplaypenalty=5000
\clubpenalty=10000
\widowpenalty=10000
\pdfinfoomitdate=1
\pdftrailerid{}
\hypersetup{pdftitle={Our definitions: Sterling constructions for e + pi + sqrt(2)}}
\newcommand{\Q}{\mathbb Q}
\newcommand{\N}{\mathbb N}
\newcommand{\G}{\mathbb Q(i)}
\newcommand{\sq}[1]{\lvert #1\rvert^2}
\newcounter{definition}
\newcommand{\entry}[2]{%
  \refstepcounter{definition}\subsection*{#1 \thedefinition: #2}}
\newcommand{\definition}[1]{\entry{Definition}{#1}}
\newcommand{\contract}[1]{\entry{Contract}{#1}}
\newcommand{\operation}[1]{\entry{Operation}{#1}}
\newcommand{\relation}[1]{\entry{Relation}{#1}}
\newcommand{\construction}[1]{\entry{Construction}{#1}}
\newcommand{\refine}{\operatorname{refine}}
\newcommand{\Poss}{\operatorname{Poss}}
\newcommand{\Seq}{\mathsf{Sequence}}
\newcommand{\interp}[1]{\mathcal I(#1)}
\newcommand{\checks}[2]{\begin{samepage}\par
  \textbf{Recognition.} #1\par\nopagebreak[4]
  \textbf{Equality.} #2\par\end{samepage}}
\newcommand{\inheritedchecks}[1]{\par
  \textbf{Recognition and equality (inherited).} #1\par}
\begin{document}
\begin{center}
  {\LARGE Our definitions\par}
  \vspace{6pt}
  {\large Sterling constructions for $e+\pi+\sqrt2$\par}
\end{center}
\vspace{8pt}

We build the finite constructions for $(e+\pi)+\sqrt2$ from boxes.
Natural numbers, integers and rationals use the box representations,
arithmetic and order in Definitions~\ref{def:naturals},
\ref{def:integers} and~\ref{def:rational-box}.
Construction equality is separate from numerical agreement of
interpretations, written $\approx$. Contracts state laws of the fixed
interpreters; operations act on the defined finite types.

\definition{Boxes and their arithmetic}\label{def:boxes}
A box is a finite expression $[B_1,\ldots,B_n]$ whose entries are
previously constructed boxes. The empty box is $0=[\,]$. Entry order
is immaterial and repetitions are retained.

The cardinality box $\#B$ replaces every immediate entry of $B$ by $0$,
retaining multiplicity.
The sum $B\uplus C$ joins the entries of the two boxes. Write $B\preceq C$
when each occurrence in $B$ can be matched to a distinct equal occurrence
in $C$. For this checked inclusion, $C\mathbin{\ominus}B$ is the box
of unmatched occurrences.

\checks{Parse balanced brackets and separators, then recursively check
each entry. Expand only previously defined finite abbreviations.
For $B\preceq C$, scan the entries of $B$, removing one equal entry from
a working copy of $C$ each time; reject if a match is absent.}
{Recursively replace each entry by its canonical bracket expression,
sort these expressions lexicographically, and enclose the sorted entries
in brackets. Compare the resulting expressions. Fix the symbol order
$[\ <\ {,}\ <\ ]$; a proper prefix precedes its extension.
Each recursion descends to smaller boxes, and each sort and scan is finite.}

\definition{Natural numbers and their arithmetic}\label{def:naturals}
A natural number is a box whose immediate entries are all $0$.
Write $n$ for the box of $n$ empty entries; in particular,
$0=[\,]$, $1=[0]$, and the successor of $n$ is $n\uplus[0]$.

\checks{Recognize a box and check that every immediate entry is empty.}
{Use box equality from Definition~\ref{def:boxes}; for natural boxes
this compares their numbers of immediate entries.}

Natural addition is box sum: $n+m=n\uplus m$. To multiply $n$ by $m$,
start with $0$ and add $m$ once for each entry of $n$.
Define $n\leq m$ by $n\preceq m$, and $n<m$ by $n\leq m$ and $n\ne m$.
When $m\leq n$, subtraction is $n-m=n\mathbin{\ominus}m$.
Powers use repeated multiplication, with $n^0=1$.

For $d>0$, divide a natural $a$ by starting with quotient $0$ and
remainder $a$. While the remainder is at least $d$, subtract $d$ and
increment the quotient. Each subtraction removes at least one entry,
so at most the original $a$ subtractions occur. The outputs satisfy
$a=qd+r$, $0\leq r<d$.
Write them as $\lfloor a/d\rfloor$ and $a\bmod d$.

Compute $\gcd(a,d)$ by the recursion
\[
 \gcd(a,0)=a,\qquad
 \gcd(a,d)=\gcd(d,a\bmod d)\quad(d>0).
\]
The positive second argument strictly decreases. Binary length is obtained
by repeatedly taking the quotient on division by $2$ until zero, counting
the steps. These procedures use only natural boxes and their finite scans.

\definition{Lists as boxes of prefixes}\label{def:lists}
For box-represented entries $a_1,\ldots,a_n$, define
\[
 \langle a_1,\ldots,a_n\rangle=[P_1,\ldots,P_n],
 \qquad P_j=[a_1,\ldots,a_j].
\]
The empty list is $[\,]$. In particular,
\[
                 \langle a,b,c\rangle=[[a],[a,b],[a,b,c]].
\]
The outer box has no prescribed entry order. Its entries are distinguished
by their cardinalities $1,\ldots,n$. Put $P_0=[\,]$; the $j$th list entry
is the unique entry of the singleton box $P_j\mathbin{\ominus}P_{j-1}$.

\checks{Recognize a box $L$ and put $n=\#L$. Find exactly one immediate
entry $P_j$ of cardinality $j$ for each $1\leq j\leq n$. Check
$P_{j-1}\preceq P_j$, remove its occurrences, and check that the
remainder has one entry. For a typed list, recognize that recovered entry
in the declared element type. Every scan is bounded by the supplied box.}
{Use box equality on canonical entry representations. Equivalently,
compare lengths and then the recovered entries at positions $1,\ldots,n$.
Repeated entries remain separate occurrences in each prefix.}

Appending $a$ to a list $L$ of length $n$ is
\[
       L\uplus[P_n\uplus[a]],\qquad P_0=[\,]\text{ when }n=0.
\]
For $n>0$, removing the first entry gives
$[P_2\mathbin{\ominus}P_1,\ldots,P_n\mathbin{\ominus}P_1]$,
with the empty case when $n=1$. A list $L$ is a prefix of a list $K$
exactly when $L\preceq K$ as boxes.
If $K=[Q_1,\ldots,Q_m]$ is another prefix-encoded list, concatenation is
\[
       L\mathbin{\Vert}K
          =L\uplus[P_n\uplus Q_1,\ldots,P_n\uplus Q_m].
\]
Indexing uses the singleton difference above. A traversal visits
$j=1,\ldots,n$; mapping, folding and bounded lookup use this finite
traversal. All ordered lists below use this encoding.
References to the first or next entry, or to traversal in order, use these
recovered positions.

\definition{Bits, pairs, records and construction trees}\label{def:finite}
A bit is the natural box $0$ or $1=[0]$. A unit has the single
representation $\star=0$. An ordered pair is
\[
                     (a,b)=\langle a,b\rangle=[[a],[a,b]].
\]
An ordered tuple is a prefix-encoded list of its fields. A record uses
the field order specified by its definition. A tagged record with tag
$t$ and fields $a_1,\ldots,a_k$ is $\langle t,a_1,\ldots,a_k\rangle$,
where $t$ is a natural box.

A construction tree is a tagged record whose prescribed fields may
themselves be construction trees. Each displayed alternative has its own
tag and field types. Within each type, assign tags $0,1,\ldots$ in the
order in which its alternatives are listed. Printed constructor names
abbreviate these tagged box records.

\checks{Decode the list by Definition~\ref{def:lists}. Check its length,
tag and prescribed field types recursively. A pair has exactly two fields;
a fixed record has exactly its declared fields. A tree child is contained
in a smaller box, so recursive recognition terminates.}
{Compare tags and corresponding fields using their declared equalities.
Use canonical box representations of fields when forming the enclosing
list. Thus tuple and tree equality preserve field order.}

A retained execution record is the tuple of its fixed algorithm tag,
inputs, list of intermediate states and output. Recognize it by
recomputing every transition and the terminal output. Its equality
compares these fields. Successful and failed branch tags are distinct.
A word is a prefix-encoded list of bits.

\definition{Integers and their arithmetic}\label{def:integers}
An integer is a signed-magnitude box pair $\langle\epsilon,m\rangle$,
where $\epsilon$ is $0$ or $1$ and $m$ is natural.
Call the codes with $\epsilon=0$ nonnegative and those with $\epsilon=1$
negative. Zero has the single
representation $\langle0,0\rangle$.
The constructor $\operatorname{int}(\epsilon,m)$ replaces either sign
at magnitude zero by this zero representation. Write the other codes
as $m$ or $-m$ according to their sign.

\checks{Decode a pair by Definition~\ref{def:finite}. Check the sign
and natural magnitude by Definition~\ref{def:naturals}, and require
$m=0\Longrightarrow\epsilon=0$.}
{Compare sign and magnitude, equivalently their canonical box pairs.}

Embed a natural $m$ as $\operatorname{int}(0,m)$.
To add equal signs, add magnitudes and retain the sign, normalizing zero.
To add opposite signs, compare magnitudes, subtract the smaller from the
larger, and use the larger magnitude's sign; equal magnitudes give zero.
Negation reverses the sign of a nonzero integer and fixes zero.
Subtraction is addition of the negative. Multiplication multiplies
magnitudes, uses sign $0$ for equal signs and $1$ for different signs,
and normalizes zero.

For order, a negative integer precedes every nonnegative integer.
Among nonnegative integers compare magnitudes in natural order;
among negative integers reverse that order.
The absolute magnitude $|a|$ is the retained natural field $m$.
All operations use the finite natural arithmetic of
Definition~\ref{def:naturals}.

\definition{Rationals and their arithmetic}\label{def:rational-box}
A rational is a reduced signed fraction encoded by
$\langle\epsilon,m,d\rangle$, where $\epsilon$ is $0$ or $1$,
$m,d$ are natural, and
\[
          d>0,\qquad \gcd(m,d)=1,\qquad
          m=0\Longrightarrow\epsilon=0.
\]
Fraction notation uses the normalization constructor below.
For an integer $a=\langle\epsilon,m\rangle$ and natural $d>0$, define
the fraction constructor by
\[
 g=\gcd(m,d),\qquad
 \operatorname{rat}(a,d)=\langle\epsilon,m/g,d/g\rangle .
\]
The two divisions here are exact natural quotients. Thus zero is
$\langle0,0,1\rangle$. Embed an integer $a$ as $\operatorname{rat}(a,1)$;
natural inputs first use their integer embedding.
Write a reduced code as $\pm m/d$ using its sign and fields.

\checks{Decode a triple using Definition~\ref{def:finite}. Check its sign
and natural fields by Definition~\ref{def:naturals}, then check the
three displayed conditions using natural order and greatest common divisor.}
{Compare the three canonical fields as boxes. For two fraction
presentations $\operatorname{rat}(a,d)$ and $\operatorname{rat}(b,e)$,
equivalently check the integer equality $ae=bd$.}

For signed integers $a,b$ and positive natural denominators $d,e$, define
\[
\begin{aligned}
 \operatorname{rat}(a,d)+\operatorname{rat}(b,e)
     &=\operatorname{rat}(ae+bd,de),\\
 -\operatorname{rat}(a,d)&=\operatorname{rat}(-a,d),\\
 \operatorname{rat}(a,d)\operatorname{rat}(b,e)
     &=\operatorname{rat}(ab,de).
\end{aligned}
\]
Products with a natural denominator use its integer embedding when
forming numerators. Subtraction adds the negative.
For $a\ne0$, the reciprocal has natural numerator magnitude $d$, the
sign of $a$, and denominator $|a|$; normalize with
$\operatorname{rat}$. Division multiplies by this guarded reciprocal.
Order is the finite cross-product test
\[
 \operatorname{rat}(a,d)\leq\operatorname{rat}(b,e)
                         \quad\Longleftrightarrow\quad ae\leq bd
\]
using integer order.

For a nonnegative rational $q=\langle0,m,d\rangle$, compute
$\lceil q\rceil$ by scanning natural $k$ from $0$ through $m$ and returning
the first with $m\leq kd$. This scan always succeeds since $d\geq1$.
The fraction readback is its retained sign, magnitude and positive denominator.
Write $\rho(q)$ for its canonical box when stored as a field below;
fraction expressions are first evaluated by these normalization rules.

\definition{Gaussian rational arithmetic}\label{def:gaussian}
A Gaussian rational is an ordered pair $z=(x,y)$ with $x,y\in\Q$, written
$x+iy$. Its box representation is $\langle\rho(x),\rho(y)\rangle$.
Equality is coordinate equality. Define
\[
\begin{aligned}
 (x,y)+(u,v)&=(x+u,y+v),&
 -(x,y)&=(-x,-y),\\
 (x,y)(u,v)&=(xu-yv,xv+yu),&
 \overline{(x,y)}&=(x,-y),\\
 \sq{(x,y)}&=x^2+y^2,&
 (x,y)^{-1}&=\frac{(x,-y)}{x^2+y^2}\quad(x^2+y^2>0).
\end{aligned}
\]
Embed a rational $q$ as $(q,0)$; put $i=(0,1)$ and $\gamma=1+i$.
Powers and finite sums use repeated multiplication and addition, with
$z^0=1$ and an empty sum equal to zero.
\checks{Accept an ordered pair of rational expressions.}
{Compare the represented rational coordinates.}

\begin{samepage}
\definition{Intervals, balls and requests}\label{def:ball}
A rational interval $[L,U]$ retains $L,U\in\Q$ and $L\leq U$.
This interval notation abbreviates the ordered box record
$\langle\rho(L),\rho(U)\rangle$.
Its width is $U-L$ and its midpoint is $(L+U)/2$.
A rational ball is a pair $(c,r)$ with $c\in\G$ and $r\in\Q$, $r\geq0$;
its membership condition for a supplied $z\in\G$ is
\[
                         \sq{z-c}\leq r.
\]
Here $r$ is a \emph{squared} radius. A squared accuracy request is a
rational $T>0$. A width request is a rational $\varepsilon>0$.
\par\end{samepage}
\checks{Check the rational fields and the displayed order or positivity
conditions by rational comparison.}
{Compare interval endpoints, or ball centre and squared radius, or request
value, as appropriate. Ball equality here means equality of these
enclosure data.}

\definition{Construction names}\label{def:codes}
The admissible construction codes are the finite expressions
\[
\begin{aligned}
 P::={}&\mathsf{Exact}(q)\mid\mathsf{HornerE}\mid\mathsf{FactorialE}
       \mid\mathsf{IntegralPi}\\
      &{}\mid\mathsf{PolygonPi}\mid\mathsf{RootTwo}
       \mid\mathsf{Add}(P,P),\qquad q\in\Q.
\end{aligned}
\]
Each tag selects the fixed algorithm defined below. The construction
object is this code with the following equality.
The grammar uses the tagged box records of Definition~\ref{def:finite};
in its displayed order, $\mathsf{Exact}$ has tag $0$ and $\mathsf{Add}$
has tag $6$. In particular,
\[
 \mathsf{Exact}(q)=\langle0,\rho(q)\rangle,\qquad
 \mathsf{HornerE}=\langle1\rangle,\qquad
 \mathsf{Add}(P,Q)=\langle6,P,Q\rangle.
\]
\checks{Check the outer tag and its exact field count. At
$\mathsf{Exact}$ invoke rational recognition; at $\mathsf{Add}$ recursively
recognize both children; the five other tags have no fields. Reject every
other expression.}
{Different outer tags are unequal. Compare rational values at
$\mathsf{Exact}$; identical parameter-free tags are equal. At
$\mathsf{Add}$ compare the left children and the right children
recursively. These recursions inspect smaller finite trees.}

Use $e_{\rm H},e_{\rm F},\pi_{\rm I},\pi_{\rm P},s$ as abbreviations for
the five parameter-free tags in their displayed order, and $P+Q$ for
$\mathsf{Add}(P,Q)$. Thus the two $e$ constructions have different
identities, as do the two $\pi$ constructions.
Write $\interp P$ for the numerical interpretation of code $P$.

\definition{States of a construction}\label{def:states}
A state is the ordered box record $\langle P,\sigma\rangle$, retaining
its construction owner $P$ and state data of this shape:
\[
\begin{array}{c|c|c}
P&\text{state data}&\text{initial data}\\ \hline
\mathsf{Exact}(q)&\star&\star\\
e_{\rm H},e_{\rm F},\pi_{\rm P},s&n\in\N&0\\
\pi_{\rm I}&\text{ordered finite list of positive width requests}&[\,]\\
\mathsf{Add}(P,Q)&(\sigma,\tau)\text{ of the two child owners}
 &(\sigma_*^P,\sigma_*^Q)
\end{array}
\]
The integral request list records continuations after its fixed initial
preparation. Frontiers, caches and intermediate states are obtained by
replaying the indicated algorithms.
\checks{Recognize the owner and recursively check its prescribed state
data, including positivity of every integral width request. Recompute
any supplied derived fields and reject a mismatch.}
{Compare owners first, then their state data recursively, retaining the
order of integral requests and of addition operands. Replayed fields
are determined by these data.}

\contract{Sample approximations}\label{def:approximation}
A sample approximation is the following numerical contract for a fixed
sampler algorithm. It supplies procedures
$n\mapsto a_n\in\G$ and $n\mapsto R_n\in\Q$, a rational bound $B\geq0$,
and certificates of
\[
 R_n\geq0,\qquad \sq{a_n}\leq B,\qquad
 \sq{a_m-a_n}\leq R_n\quad(m\geq n).
\]
Its refinement procedure, on a current stage $c\in\N$ and $T>0$,
returns an actual $j\geq c$ with $R_j<T$, together with these two
checkable inequalities.
The sampler-code family here is precisely $e_{\rm H},e_{\rm F},
\pi_{\rm P},s$, with the algorithms specified below. The universal laws
are correctness obligations for these fixed algorithms.
\checks{For this family, use Definition~\ref{def:codes} and accept only
the four parameter-free tags $\mathsf{HornerE}$, $\mathsf{FactorialE}$,
$\mathsf{PolygonPi}$ and $\mathsf{RootTwo}$.}
{Use Definition~\ref{def:codes}: two sampler codes are equal exactly
when their tags agree. These checks concern the codes; the displayed
universal laws are the interpretation contract.}

\definition{Geometric selector codes}\label{def:selector}
For $C\in\Q$, $C\geq0$, define $\mathsf{Select}_C(c,T)$ for $c\in\N$
and $T>0$ as follows. Compute
\[
 h=\lceil C/T\rceil,\qquad
 f=\begin{cases}
 1,&h=0,\\
 \text{number of binary digits of }h,&h>0.
 \end{cases}
\]
Start at stage $j=c$ with a natural counter of value $f$. If $C/4^j<T$,
return $j$. Otherwise, if the counter is positive, increment $j$ and
decrement the counter; repeat. Return $c+f$ when the counter reaches zero.
The selector retains the input, returned stage and certificate
\[
                 c\leq j\leq c+f,\qquad C/4^j<T.
\]
The counter is the procedure's \emph{fuel}: an explicit bound on the
number of stage advances.
The selector code is $\mathsf{Select}(C)$.
\checks{Check this tag, one rational field $C$, and $C\geq0$.
An invocation additionally requires $c\in\N$ and $T>0$.}
{Compare the rational parameter $C$; compare all retained input and output
fields when comparing invocation records.}

\contract{Persistent numerical processes}\label{def:process}
The numerical interpretation of a construction code has the following
operations and laws on its states from Definition~\ref{def:states}:
\begin{itemize}
\item A ball $(c_s,r_s)$ at each state $s$, and a possible-reading
relation $\Poss(s,z)$ for $z\in\G$. Evidence for a possible reading is
retained with that reading.
\item A witness $w_s\in\G$ with evidence of $\Poss(s,w_s)$, and a
containment law
\[
                    \Poss(s,z)\ \Longrightarrow\ \sq{z-c_s}\leq r_s.
\]
\item A refinement $s'=\refine(s,T)$ for every $T>0$, with laws
\[
 r_{s'}<T,\qquad \Poss(s',z)\ \Longrightarrow\ \Poss(s,z).
\]
\item A rational bound $B\geq0$ and the law $\sq{c_s}\leq B$ at every
state.
\end{itemize}
These laws govern the fixed interpretation procedures below.
\inheritedchecks{Construction owners use Definition~\ref{def:codes},
states use Definition~\ref{def:states}, and observed balls use
Definition~\ref{def:ball}. The contract states laws of their interpretation;
it supplies no additional type of arbitrary function-bearing records.}

\contract{Coherence}\label{def:name}
A coherent numerical interpretation supplies a rational allowance
$\eta(s)\geq0$ at every state, with
\[
\begin{gathered}
 r_s\leq\eta(s),\\
 \Poss(s,u),\ \Poss(t,v)
 \ \Longrightarrow\
 \sq{u-v}\leq4\bigl(\eta(s)+\eta(t)\bigr).
\end{gathered}
\]
It also supplies a request modulus $M(A)>0$ for every rational $A>0$,
with the law
\[
             0<r<M(A)\ \Longrightarrow\
             \eta\bigl(\refine(s,r)\bigr)<A.
\]
A \emph{name} uses a construction code with the initial state prescribed by
Definition~\ref{def:states}. Its numerical interpretation has the displayed
coherence laws. The finite code determines the algorithms, allowances and
modulus; these are not arbitrary functions supplied for recognition.
\inheritedchecks{Use Definition~\ref{def:codes} for the name's code and
Definition~\ref{def:states} for its prescribed initial state. Allowances
and moduli are computed by the fixed interpreter; the coherence laws
are its contract.}

\operation{Sequence-backed interpretation}\label{def:sequence}
For one of the four sampler codes, write $\Seq(a,R,B)$ for its
sequence-backed numerical interpretation. It has
states $n\in\N$, initially $0$, and
\[
 (c_n,r_n)=(a_n,R_n),\qquad w_n=a_n,\qquad
 \eta(n)=R_n,\qquad M(A)=A.
\]
A possible-reading certificate at $n$ retains a natural $m\geq n$
and the exact equality $z=a_m$. Refinement is the approximation's
refinement from the current stage $n$. Its centre bound is $B$.
\checks{For the sampler owner, use the four-code restriction in
Contract~\ref{def:approximation}; for its state use
Definition~\ref{def:states}. A supplied possible-reading witness must have
a natural index $m\geq n$ and reading equal to the recomputed $a_m$.}
{Owner and state equality are inherited from
Definitions~\ref{def:codes} and~\ref{def:states}.
Compare witness indices as natural boxes and readings by
Definition~\ref{def:gaussian}, using the record rule of
Definition~\ref{def:finite}.}

\operation{Canonical sampling}\label{def:sampling}
For a name with initial state $s_*$, define its canonical states and
samples by the finite recursion
\[
 s_0=\refine(s_*,1),\qquad
 s_{k+1}=\refine(s_k,1/4^{k+1}),\qquad A_k=w_{s_k}.
\]
A canonical stage counts this recursion; it need not equal a raw stage
of the name's original sample procedure.
\inheritedchecks{Input codes use Definition~\ref{def:codes} and stage
indices use natural boxes from Definition~\ref{def:naturals}. Output states
use Definition~\ref{def:states} and samples use
Definition~\ref{def:gaussian}. Check any retained execution record by
replaying the recursion, as in Definition~\ref{def:finite}.}

\relation{Operational equivalence}\label{def:equivalence}
For names $A,B$, a comparison at independent canonical floors $p,q\in\N$
and squared request $T>0$ retains both names and supplies $i\geq p$,
$j\geq q$, with a finite certificate of
\[
              \sq{A_m-B_n}<T
              \qquad\text{for every }m\geq i,\ n\geq j.
\]
Write $A\approx B$ when a finite algorithm with a finite correctness and
termination argument constructs such a comparison
for every $p,q,T$. This relation is distinct from presentation identity.
It concerns the numerical interpretations; the sterling equality of
construction objects remains Definition~\ref{def:codes}. Finite comparison
derivations are defined separately below.
The arguments are recognized construction names. This entry specifies a
relation between their numerical interpretations; it supplies no new
object type whose equality is $\approx$.

\definition{Finite comparison derivations}\label{def:derivations}
A comparison derivation is a finite tree with the following tags and
endpoint rules. Write $D:P\approx Q$ for its computed endpoints.
Each tree node uses a tagged prefix-encoded record.
\[
\begin{array}{l|l}
\text{tag and fields}&\text{computed endpoints}\\ \hline
\mathsf{Refl}(P)&P\approx P\\
\mathsf{SameE}&e_{\rm H}\approx e_{\rm F}\\
\mathsf{SamePi}&\pi_{\rm I}\approx\pi_{\rm P}\\
\mathsf{Sym}(D),\quad D:P\approx Q&Q\approx P\\
\mathsf{Trans}(D,F),\quad D:P\approx Q,\ F:Q\approx R&P\approx R\\
\mathsf{Sum}(D,F),\quad D:P\approx Q,\ F:R\approx S
 &(P+R)\approx(Q+S)\\
\mathsf{Comm}(P,Q)&(P+Q)\approx(Q+P)\\
\mathsf{Assoc}(P,Q,R)&((P+Q)+R)\approx(P+(Q+R))\\
\mathsf{Zero}(P)&(P+\mathsf{Exact}(0))\approx P\\
\mathsf{RationalSum}(q,r)&
 (\mathsf{Exact}(q)+\mathsf{Exact}(r))\approx\mathsf{Exact}(q+r)
\end{array}
\]
The table fixes the tag order. Thus $\mathsf{SameE}=\langle1\rangle$
and $\mathsf{Sum}(D,F)=\langle5,D,F\rangle$.
Each tag selects the corresponding fixed comparison construction:
$\mathsf{SameE}$ uses the Horner--factorial comparison;
$\mathsf{SamePi}$ uses the integral--polygon comparison; the other tags
use the indicated reflexive, symmetric, transitive or additive comparison.
The comparison arguments are given in
\href{SterlingArithmetic.pdf}{\emph{Sterling arithmetic}},
Sections~3, 5--7 and 9.

\checks{Traverse the tree, check tags and field types, and compute its
endpoints. At $\mathsf{Trans}$, check equality of the two middle
construction codes by Definition~\ref{def:codes}. Reject every unlisted
tag or mismatched endpoint.}
{Compare derivation tags and fields recursively. Distinct derivation trees
with the same endpoints remain distinct derivation objects.}

A checked derivation selects an algorithm satisfying
Relation~\ref{def:equivalence}. The recognition test concerns the supplied
tree; it is not a search for some derivation with given endpoints.

\operation{Exact rational interpretation}\label{def:exact}
For $q\in\Q$, the exact name $\mathsf{Exact}(q)$ has one state, ball
$((q,0),0)$, witness $(q,0)$ and possible-reading condition $z=(q,0)$.
Refinement returns the same state. Its centre bound is $q^2$, its
coherence allowance is zero, and its modulus is $M(A)=A$.
\inheritedchecks{Rational inputs use Definition~\ref{def:rational-box}.
The resulting $\mathsf{Exact}(q)$ code uses Definition~\ref{def:codes};
its unit state uses Definition~\ref{def:states}. Its ball and reading
use Definitions~\ref{def:ball} and~\ref{def:gaussian}.}

\operation{Addition of names}\label{def:add}
For names $A,B$, the name $A+B$ retains an ordered pair $(s,t)$ of
their original states, initially $(s_*^A,s_*^B)$. Define
\[
\begin{aligned}
 c_{(s,t)}&=c_s^A+c_t^B,&
 r_{(s,t)}&=2(r_s^A+r_t^B),\\
 w_{(s,t)}&=w_s^A+w_t^B,&
 \eta_+(s,t)&=2(\eta_A(s)+\eta_B(t)).
\end{aligned}
\]
A possible sum reading retains $u,v$, their original possible-reading
certificates, and $z=u+v$. Define
\[
 \refine_+((s,t),T)
   =\bigl(\refine_A(s,T/4),\refine_B(t,T/4)\bigr).
\]
The centre bound is $2(B_A+B_B)$. For positive $u,v$, put
$\operatorname{com}(u,v)=uv/(u+v)$; the sum modulus at $a>0$ is
\[
 M_+(a)=4\operatorname{com}\bigl(M_A(a/4),M_B(a/4)\bigr).
\]
\inheritedchecks{Recognize both operand codes and the resulting
$\mathsf{Add}$ code by Definition~\ref{def:codes}. Their states and the
ordered pair state use Definition~\ref{def:states}; balls use
Definition~\ref{def:ball}. Sum-code equality compares the two ordered
child codes, as in Definition~\ref{def:codes}.}

\operation{Horner's description of \texorpdfstring{$e$}{e}}\label{def:horner}
For $z\in\G$, positive natural $d$, and $k,n\in\N$, define
\[
\begin{aligned}
 H_z(d,0)&=1,&
 H_z(d,k+1)&=1+\frac{zH_z(d+1,k)}{d},\\
 E_n(z)&=1+zH_z(2,n),&
 h_n&=E_n(1).
\end{aligned}
\]
The numerical interpretation of $\mathsf{HornerE}$ is
\[
             \interp{e_{\rm H}}=\Seq\bigl(h,\ (10/4^n)_{n},\ 10\bigr),
\]
with refinement $\mathsf{Select}_{10}$. The procedure $E_n(1)$ is the
Horner exponential procedure evaluated at one.
\inheritedchecks{The constant's code is the parameter-free
$\mathsf{HornerE}$ tag of Definition~\ref{def:codes}. Its stage is a
natural box (Definition~\ref{def:naturals}); its state and sample use
Definitions~\ref{def:states} and~\ref{def:gaussian}.
For the helper $H_z(d,k)$, additionally recognize $z$ by
Definition~\ref{def:gaussian} and check natural $d>0$ and $k$.}

\operation{The factorial-sum description of \texorpdfstring{$e$}{e}}\label{def:factorial}
Define $0!=1$, $(k+1)!=(k+1)k!$, and the finite sum
\[
                         f_n=\sum_{k=0}^{n+1}\frac1{k!}.
\]
The numerical interpretation of $\mathsf{FactorialE}$ is
\[
             \interp{e_{\rm F}}=\Seq\bigl(f,\ (10/4^n)_{n},\ 10\bigr),
\]
with refinement $\mathsf{Select}_{10}$. The two names retain their distinct
procedure tags.
\inheritedchecks{The code is the parameter-free $\mathsf{FactorialE}$
tag of Definition~\ref{def:codes}. Stages use natural boxes from
Definition~\ref{def:naturals}; states use Definition~\ref{def:states}
and samples use Definition~\ref{def:gaussian}.}

\newpage
\definition{Rational source frontiers}\label{def:frontier}
Define the rational source
\[
                         g(t)=\frac4{1+t^2}\qquad(t\in\Q).
\]
A source frontier record $F$ retains a finite rational partition
$0=x_0<\cdots<x_N=1$. For each cell $[x_j,x_{j+1}]$ it retains:
its width $w_j=x_{j+1}-x_j$; a rational tag in the cell; the exact
source reading at that tag; and a rational range $[\ell_j,u_j]$
with the endpoint checks
\[
              \ell_j\leq g(x_{j+1}),\qquad
              g(x_j)\leq u_j.
\]
Its source-range law is
\[
              \ell_j\leq g(t)\leq u_j
              \qquad(x_j\leq t\leq x_{j+1},\ t\in\Q).
\]
Its enclosure is
\[
             I(F)=\left[\sum_{j=0}^{N-1}w_j\ell_j,\
                        \sum_{j=0}^{N-1}w_ju_j\right].
\]
Its box record is
\[
\begin{aligned}
 C_j&=\langle x_j,x_{j+1},w_j,t_j,g(t_j),\ell_j,u_j\rangle,\\
 F&=\bigl\langle
       \langle x_0,\ldots,x_N\rangle,\
       \langle C_0,\ldots,C_{N-1}\rangle
     \bigr\rangle ,
\end{aligned}
\]
where every rational field uses $\rho$, and $t_j$ is the cell's tag.
A source cache is an ordered finite list of retained tags and exact
readings, encoded as $\langle(t_1,y_1),\ldots,(t_m,y_m)\rangle$
with $y_j=g(t_j)$. A lookup scans this list in order and reuses the first matching
reading; on a miss it evaluates $g$ and appends the new tag and reading.
\checks{Check the finite partition, ordered cuts and endpoints; one cell
per adjacent interval; each tag's membership; the exact width and source
readback; and the two rational endpoint inequalities for each range.
For a cache, check every saved reading against $g$ at its retained tag.}
{Compare the partition and ordered cell fields, including tags, readings
and ranges. Compare caches in their retained list order. An enclosure
alone does not determine frontier identity.}

\definition{Frontier preparation and refinement records}\label{def:common}
For a width request $\varepsilon>0$, define
\[
             v=\varepsilon/4,\qquad
             N=\lceil32/\varepsilon\rceil+1.
\]
The fresh frontier $F_\varepsilon$ has cuts $j/N$, $0\leq j\leq N$.
Cell $j$ has left tag $j/N$, its cached or fresh reading $g(j/N)$,
and range $[g(j/N)-v,g(j/N)+v]$. Its retained certificates include the
cell range laws and $\operatorname{width} I(F_\varepsilon)<\varepsilon$.
Process the cells from left to right, starting with the supplied cache,
which is empty for an initial preparation.

For two frontiers of $g$, their common frontier is defined by merging
their ordered cuts and removing duplicates. Each child retains its actual
parent cell from each frontier and the endpoint inequalities witnessing
both inclusions. Its range is
\[
                  [\max(\ell,\ell'),\ \min(u,u')],
\]
where the two parent ranges are $[\ell,u]$ and $[\ell',u']$.
Its tag is its left endpoint. Retain the reading there, the two parent
range certificates and the source/tag equalities for any reused reading.
Process the merged children from left to right.

A fresh-preparation record retains its positive request, input cache and
computed frontier and cache. A common-refinement record retains both
parent frontiers, its input cache, the ordered children with their two
parent indices, and the computed frontier and cache.
\checks{Recognize the inputs and recompute the specified finite
preparation or cut merge, every parent lookup and every derived field.
Compare the claimed outputs with those computed.}
{Compare the retained inputs, parents and ordered finite output fields.
The parent frontiers remain part of a refinement record's identity.}

\operation{The rational-integral description of \texorpdfstring{$\pi$}{pi}}\label{def:integral}
The empty request list is interpreted as the fresh frontier $F_1$, its
enclosure $[L_0,U_0]$ and its source cache. Interpret a nonempty state list
by replaying that preparation and then each width request in list order.

For a width request $\varepsilon$, reuse the current frontier if its
width is already below $\varepsilon$. Otherwise build $F_\varepsilon$,
using saved readings, and take its common frontier with the current one.
Retain both parents, the common frontier and the readings in the continued
history. In either case the returned enclosure has width below
$\varepsilon$. Append the request to the state list, including in the reuse
case, and retain its replayed output.

For current enclosure $[L,U]$ and $D=U-L$, define
\[
\begin{gathered}
 c=(L,0),\qquad r=D^2,\qquad
 w=((L+U)/2,0),\qquad \eta=D^2,\\
 \Poss(s,z)\ \Longleftrightarrow\
       z=(x,0)\text{ for a retained rational }L\leq x\leq U.
\end{gathered}
\]
Refinement at squared request $T>0$ is the history continuation at width
request $\varepsilon=T/(T+1)$. The centre bound is
$2(L_0^2+(U_0-L_0)^2)$ and the coherence modulus is $M(A)=A$.
These operations define $\interp{\pi_{\rm I}}$.
\inheritedchecks{The code is the parameter-free $\mathsf{IntegralPi}$
tag of Definition~\ref{def:codes}. States use Definition~\ref{def:states},
including the whole ordered request list; requests use
Definition~\ref{def:ball}. Replayed frontiers and caches use
Definition~\ref{def:frontier}, and continuation records use
Definition~\ref{def:common}.}

\operation{The rational-polygon description of \texorpdfstring{$\pi$}{pi}}\label{def:polygon}
For rational $t\in[0,1]$, define
\[
 p(t)=\left(\frac{1-t^2}{1+t^2},\frac{2t}{1+t^2}\right),
 \qquad R(x,y)=(-y,x),\qquad
 \det((x,y),(u,v))=xv-yu.
\]
For positive natural $N$, form the ordered list of $4N$ vertices
\[
             v_{qN+j}=R^q p(j/N)
             \quad(0\leq q<4,\ 0\leq j<N).
\]
Close the list by $v_{4N}=v_0$ and define its rational area
\[
                   A_N=\frac12\sum_{j=0}^{4N-1}\det(v_j,v_{j+1}).
\]
The numerical interpretation of $\mathsf{PolygonPi}$ is
\[
           \interp{\pi_{\rm P}}=
           \Seq\bigl((A_{2^n},0)_{n},\ (4/4^n)_{n},\ 16\bigr),
\]
with refinement $\mathsf{Select}_4$.
\inheritedchecks{The code is the parameter-free $\mathsf{PolygonPi}$
tag of Definition~\ref{def:codes}. Its natural stage and positive cell
count use Definition~\ref{def:naturals}; its state and sample use
Definitions~\ref{def:states} and~\ref{def:gaussian}.
The ordered vertex list uses Definition~\ref{def:lists}.}

\operation{Positive-root sampling}\label{def:root}
Define the natural square-root procedure $\operatorname{isqrt}$ by
$\operatorname{isqrt}(0)=0$. For $x>0$, set
\[
 k=\operatorname{isqrt}(\lfloor x/4\rfloor),\qquad b=2k,
 \qquad
 \operatorname{isqrt}(x)=
 \begin{cases}
 b+1,&(b+1)^2\leq x,\\
 b,&(b+1)^2>x.
 \end{cases}
\]
It returns $m$ with the certificate $m^2\leq x<(m+1)^2$.
For stage $n$, define
\[
 d_n=2^n,\qquad m_n=\operatorname{isqrt}(2d_n^2),\qquad
 L_n=\frac{m_n}{d_n},\qquad U_n=\frac{m_n+1}{d_n}.
\]
The retained bracket certificate is
\[
          0\leq L_n,\qquad L_n^2\leq2<U_n^2,\qquad
          U_n-L_n=1/2^n.
\]
The numerical interpretation of $\mathsf{RootTwo}$ is
\[
             \interp{s}=\Seq\bigl((U_n,0)_{n},\ (1/4^n)_{n},\ 9\bigr),
\]
with refinement $\mathsf{Select}_1$. The notation $\sqrt2$ denotes this
name.
\inheritedchecks{The code is the parameter-free $\mathsf{RootTwo}$
tag of Definition~\ref{def:codes}. Its stage and the input to
$\operatorname{isqrt}$ use natural boxes (Definition~\ref{def:naturals});
states use Definition~\ref{def:states}. Bracket intervals and samples
use Definitions~\ref{def:ball} and~\ref{def:gaussian}, respectively.
Recompute any retained root trace by Definition~\ref{def:finite}.}

\definition{Projective readings and decoder matrices}\label{def:projective}
A projective reading is presented by a pair $[u:v]$ of Gaussian
rationals with at least one coordinate nonzero. Its equality is
\[
                         [u:v]=[u':v']
                  \quad\Longleftrightarrow\quad uv'=u'v.
\]
An affine reading $z$ is presented by $[z:1]$; $[1:0]$ is the terminal
tail tag. Its canonical representative for storage as a box field is
\[
 \nu([u:v])=
 \begin{cases}
   \langle u/v,1\rangle,&v\ne0,\\
   \langle1,0\rangle,&v=0.
 \end{cases}
\]
The entries use Gaussian rational box representations.
A decoder matrix is an ordered $2$ by $2$ table of Gaussian
rationals with nonzero determinant. It acts by
\[
 \begin{pmatrix}a&b\\c&d\end{pmatrix}[u:v]=[au+bv:cu+dv].
\]
Its box representation is
$\langle\langle a,b\rangle,\langle c,d\rangle\rangle$.
Matrix multiplication uses the two-term row-by-column sums.
\checks{For a projective presentation, check the two Gaussian fields
and that they are not both zero. For a decoder matrix, check four
Gaussian fields and $ad-bc\ne0$.}
{Compare projective readings by the displayed cross-product test.
Equivalently, compare their canonical representatives $\nu$ as boxes.
Compare decoder matrices entry by entry. Matrices with equal projective
actions may still be distinct matrix objects.}

\definition{CCL blocks and their readings}\label{def:blocks}
A raw complex continued logarithm (CCL) word is a finite ordered list of
bits. Displayed bit strings abbreviate prefix-encoded lists; juxtaposition
means their concatenation $\Vert$. The notation for a completed ordinary
block abbreviates the word
\[
                    [k,a]=1^k0\,\operatorname{code}(a).
\]
The initial ordinary block has $k\geq1$; later ordinary blocks have
$k\geq0$. A low initial block is $0[k,a]$, with $k\geq0$.
For a normalized reading $w=x+iy$, the sector table is
\begingroup
\renewcommand{\arraystretch}{1.25}
\[
\begin{array}{c|c|c|c}
x&y&a&\operatorname{code}(a)\\ \hline
x\geq\tfrac12&-\tfrac12<y<\tfrac12&1&\texttt{000}\\
x\geq\tfrac12&y\geq\tfrac12&1+i&\texttt{001}\\
-\tfrac12<x<\tfrac12&y\geq\tfrac12&i&\texttt{010}\\
x\leq-\tfrac12&y\geq\tfrac12&-1+i&\texttt{011}\\
x\leq-\tfrac12&-\tfrac12<y<\tfrac12&-1&\texttt{100}\\
x\leq-\tfrac12&y\leq-\tfrac12&-1-i&\texttt{101}\\
-\tfrac12<x<\tfrac12&y\leq-\tfrac12&-i&\texttt{110}\\
x\geq\tfrac12&y\leq-\tfrac12&1-i&\texttt{111}
\end{array}
\]
\endgroup
The canonical normalization predicates for a supplied reading $z$ use
$w=z/\gamma^k$ in an ordinary block and $w=\gamma^kz$ in a low initial
block. They require $\sq w\in[1,2)$ initially and $\sq w\in(1,2]$
later, together with the selected row of the sector table.
Each predicate is a finite test on the supplied rational coordinates.

For a supplied next tail $t$, the ordinary and low block decoders are,
respectively,
\[
        \gamma^k\left(a+\frac1{\gamma t}\right),
        \qquad
        \gamma^{-k}\left(a+\frac1{\gamma t}\right).
\]
Their homogeneous matrices are
\[
 B_{\rm ordinary}=
 \begin{pmatrix}\gamma^{k+1}a&\gamma^k\\\gamma&0\end{pmatrix},
 \qquad
 B_{\rm low}=
 \begin{pmatrix}\gamma a&1\\\gamma^{k+1}&0\end{pmatrix}.
\]
A terminating tail is $[1:0]$. Consecutive block decoders compose by matrix
multiplication in their word order, using Definition~\ref{def:projective}.
\checks{For a raw word, recognize a finite bit list. For a completed block,
check its kind, unary count, separator and three sector bits; the initial
ordinary count must be positive.}
{Compare raw words bit by bit. Compare block kind, natural count and
sector code. Projective readings of supplied tails use the separate
cross-product equality.}

\definition{CCL parser records}\label{def:ccl-parser}
A parser record retains a raw word $w$, its current phase $\varphi$ and
the decoder matrix $M$ of its completed blocks:
\[
                         \langle w,\varphi,M\rangle.
\]
The phase retains its kind, unary count and any unfinished sector bits.
The empty word has initial phase and identity matrix.

Parse the bits in their recovered list order. A first bit of one selects
initial-high count one; a first bit of zero selects initial-low count zero.
In a unary phase, a one increments the count and a zero starts a three-bit
sector. On completing a sector, multiply $M$ on the right by that block's
decoder matrix from Definition~\ref{def:blocks}, then begin the later unary
phase at count zero. Retain incomplete unary and sector phases.

\checks{Recognize the raw word by Definition~\ref{def:blocks} and replay
the parser fold. Check the supplied phase, count, partial sector bits and
matrix against the computed fields.}
{Compare the retained word and phase fields by their finite-data
equalities, and the decoder matrix by Definition~\ref{def:projective}.}

\definition{CCL cursor records}\label{def:ccl-cursor}
The CCL output algorithm is the fixed four-mode Purelib transducer:
\texttt{Machine.initial} starts it, \texttt{Machine.Step.nextBit} returns
one bit, and \texttt{Machine.Advance.advance} iterates these steps.
Its mode-branch procedure is
\texttt{Emission.Step.step}.\footnote{Under
\path{Purelib/Data/CCL/Transducer/}, the source definitions are in
\path{Machine/Basic.lean}, \path{Machine/Step.lean},
\path{Machine/Advance.lean} and \path{Emission/Step.lean}.}
The box records below encode its finite execution data. Its rational
branch tests, process charts and nonempty emitted chunks are an explicit
dependency here.

A cursor record retains its construction origin, a mode, its current
numerical state, an emitted word, a pending word and its finite history.
These are its ordered box-record fields. Both words and the history use
prefix-encoded lists. The origin is a retained construction owner and
state, at which the cursor starts with empty emitted and pending words.
Its \emph{committed word} is emitted followed by pending.

The cursor modes are initial, low, loop and residual, each retaining the
finite process-chart code and encoded state generated by the fixed cursor.
Chart construction and state replay use that algorithm's finite records.
Low and loop also retain their scale count;
loop retains whether it is initial-high or later. Residual retains a
unary count. For a possible reading $z$ of the current mode and count $k$,
the associated projective tails are
\[
\begin{array}{c|c|c|c}
\text{initial}&\text{low}&\text{loop}&\text{residual}\\ \hline
[z:1]&[\gamma^{-k}z:1]&[\gamma^kz:1]&[1:\gamma z].
\end{array}
\]

\checks{Recognize the origin by Definition~\ref{def:states} and the
words by Definition~\ref{def:blocks}. Replay the finite cursor history
with the fixed transducer; check its generated chart code, state, mode,
counts, rational bounds and emitted and pending words.}
{Compare the origin and ordered history, then the retained finite mode,
state, bound and word fields. Parser data derived from the committed word
use Definition~\ref{def:ccl-parser}.}

\operation{Advancing a CCL cursor}\label{def:ccl-advance}
To advance a cursor by one bit, use the following two cases:
\begin{enumerate}
\item If pending is nonempty, remove its first bit and append that bit
to emitted. Keep the mode and its numerical state.
\item If pending is empty, compute one nonempty chunk and the next mode
with the fixed transducer. Append the first chunk bit to emitted and
retain the remaining chunk as pending.
\end{enumerate}
Append this transition to the cursor history. To advance by a natural
count $k$, repeat this procedure $k$ times; a count of zero returns the
original cursor.

\inheritedchecks{Input and output cursors use
Definition~\ref{def:ccl-cursor}; the count uses
Definition~\ref{def:naturals}. Check a supplied transition or finite run
by replay, using the execution-record rule of Definition~\ref{def:finite}.}

\contract{CCL cursor reading and productivity}\label{def:ccl-reading}
Every advance of Operation~\ref{def:ccl-advance} terminates and appends
exactly one emitted bit. A $k$-step advance therefore increases the emitted
length by $k$ and preserves the previous emitted word as a prefix.
For the prefix-encoded words, this extension is box inclusion as in
Definition~\ref{def:lists}.

The parser phase of the committed word agrees with the retained mode:
initial mode has initial phase; low mode has low unary phase; loop mode
has the unary phase of its retained kind; residual mode has later unary
phase. The unary counts agree in each case.

The \emph{exact reading law} concerns the same retained origin.
Let $M$ be the decoder matrix of the committed word. For every supplied
possible reading of the current mode, applying $M$ to its projective tail
from Definition~\ref{def:ccl-cursor} yields $[u:1]$ for a retained reading
$u$ possible at the original state where the cursor started.
Digit decisions may use certified redundant margins and require no
decision that a remainder is zero.

\inheritedchecks{Cursor records use Definition~\ref{def:ccl-cursor},
their parser records use Definition~\ref{def:ccl-parser}, and projective
readings use Definition~\ref{def:projective}. These are laws of the fixed
transducer on those records.}

\definition{CCL accuracy observations}\label{def:ccl}
An observation retains a construction owner, its initial state, a positive
squared request $T$, a natural bit count $k$, a refined state, a rational
ball and a final cursor.

For a squared request $T>0$ and bit count $k\in\N$, a finite CCL
observation of a name starts by retaining
\[
                  s'=\refine(s_*,T),\qquad
                  (c_{s'},r_{s'}),\qquad r_{s'}<T.
\]
Start a cursor of Definition~\ref{def:ccl-cursor} at this same $s'$ with
both words empty, then advance $k$ times by
Operation~\ref{def:ccl-advance}. Return the emitted word, of length $k$,
the retained ball as its \emph{accuracy sidecar}, and the exact reading
certificate of Contract~\ref{def:ccl-reading}.

A new accuracy request may start a new cursor at a new refined state.
Prefix preservation concerns continuation of one cursor. No specified
accuracy is inferred from the bit count alone.
\checks{Check the construction owner, original state, positive request,
bit count and finite record fields. Recompute refinement and its ball.
Recognize the final cursor by Definition~\ref{def:ccl-cursor}, checking
that it starts at this refined state and records exactly $k$ advances.}
{Compare the owner, original state, request, bit count and ordered
refinement data using their previously defined equalities, and the final
cursor by Definition~\ref{def:ccl-cursor}. Equal exposed words alone do
not make two observations equal.}

\construction{The presented expression}\label{def:expression}
An expression for this example is a construction code of
Definition~\ref{def:codes}. Interpret its leaves by their definitions above
and each ordered addition node by Operation~\ref{def:add}. Numerical
agreement is the operational equivalence of Relation~\ref{def:equivalence}.

The two presentations of the requested example are
\[
                  X=(e_{\rm H}+\pi_{\rm I})+s,\qquad
                  Y=(e_{\rm F}+\pi_{\rm P})+s.
\]
Their initial states retain the corresponding ordered triples of original
leaf states, grouped as $((s_e,s_\pi),s_{\sqrt2})$. Their CCL observations
are the words, accuracy sidecars and reading certificates of
Definition~\ref{def:ccl} for these interpreted names.
The supplied comparison derivation for these two presentations is
\[
 \mathsf{Sum}\bigl(
   \mathsf{Sum}(\mathsf{SameE},\mathsf{SamePi}),\
   \mathsf{Refl}(s)\bigr).
\]
Its endpoints are $X\approx Y$.
\checks{Recognize $X$ and $Y$ by Definition~\ref{def:codes}, checking
their indicated leaf tags and ordered addition nodes. Recognize the
displayed comparison derivation by Definition~\ref{def:derivations}.}
{Use Definition~\ref{def:codes} for the two construction codes and
Definition~\ref{def:derivations} for derivation objects. The supplied
comparison establishes the separate relation $X\approx Y$.}

\end{document}
